\section{Introduction}
\label{sec:intro}
World models offer a direct route from visual prediction to control: predict the consequences of candidate actions, choose a promising sequence, and replan after execution. Joint-embedding predictive architectures make this approach practical by modeling dynamics in a compact latent space~\cite{zhou2025dinowm,assran2025vjepa2,maes2026lewm}. However, a controller can keep predicting and replanning while repeatedly returning to the same region. The current observation and its distance to the goal describe the task that remains; the recent execution history also reveals where the controller has already tried to make progress.

This distinction matters in both manipulation and navigation. A pusher may move around an almost-aligned object without reaching the required agent position. A navigator may repeatedly traverse the same corridor without reaching its destination. Such behavior supplies actionable feedback: it identifies configurations where the current combination of model, objective, and optimizer is unproductive. Our central idea is to turn that feedback into a change in the next planning objective. \Cref{fig:teaser} shows how this intervention redirects failed executions with both frozen and adaptive world models.

We introduce \emph{execution-conditioned planning} (ECP), which combines two mechanisms. An episode-local \emph{stagnation regularizer} adds a repulsive cost around observations associated with repeated local motion. Temporary proprioceptive reweighting increases the emphasis on agent positioning during recovery. The regularizer remembers \emph{where} execution stalled; reweighting specifies \emph{which goal component} should guide the correction. Both act through the existing action optimizer. The world model can therefore remain frozen or continue adapting from newly observed transitions.

The memory cost has compact support, with its radius scaled by the remaining distance to the goal. It redirects plans near a stalled configuration while leaving the goal unpenalized. This provides a planning-stage complement to learning the predictive representation: Temporal Straightening regularizes latent trajectory geometry during training~\cite{wang2026straightening}, and AdaJEPA updates model parameters during deployment~\cite{wang2026adajepa}. ECP makes action selection responsive to the behavior already observed. For adaptive control, a prediction-error gate can coordinate the onset of correction with ongoing model updates.

Training data provide a second, complementary route to stronger control. Following AdaJEPA's contact-biased collection construction~\cite{wang2026adajepa}, we augment source controls with an object-directed bias and stochastic variation, then record the interactions those controls produce. This enriches the experience used to train the original predictive architecture. The two interventions have distinct roles: data augmentation changes what the world model learns from, while ECP changes how the controller responds when execution stalls.

On PushObj, data augmentation raises success from 43.7\% to 61.1\% for frozen planning and from 69.4\% to 74.9\% with AdaJEPA. ECP increases these rates to 72.3\% and 76.9\%, respectively, under the same 30-replan budget (\cref{fig:headline}). With the augmented-data frozen model, conditional correction recovers 39 additional tasks while retaining all 214 native successes. On Maze, memory regularization improves frozen success from 48\% to 50\%; gated memory improves adaptive success from 63\% to 66\%, also preserving every native success. Controlled ablations identify the contributions of memory, goal weighting, and activation timing.

Our contributions are:
\begin{itemize}[leftmargin=*,itemsep=1pt,topsep=2pt]
\item An execution-conditioned objective that uses stalled observations to regularize predicted trajectories and temporarily reweight goal components.
\item A training and planning approach that combines interaction-biased data augmentation with ECP to improve frozen and adaptive world-model control.
\item Paired evaluations, component ablations, and verified simulator trajectories that connect improved task completion to stagnation-aware intervention.
\end{itemize}
